\section{怎样写读后感、观后感}

我们读了一篇文章或者一本著作，看了一部影片或者欣赏了一幅图画之后，总是情不自禁地说说自己读后或观后的感想吧。如果你把这种感想写成了文章，那就是人们常说的“读后感”或者“观后感”了。

读后感或者观后感，由于都是以阐发感想、写出看法为主要内容，所以都属于议论文的范畴，也是常见的作文训练方式之一。这种文章的特点是：必须从所读的文章、所看的影片、戏剧、图画出发，抓住原文或原画等的主要观点或精神实质，从而引出自己的感想来。主要感受即中心论点，必须紧扣原文，原画等生发出来，而不可与原文、原画等关系不大甚至无关；在写作方式上可以叙事、抒情和夹叙夹议，但一定要以阐发感想的议论为主。至于为了证明主要感受的正确性，要引用事例或引用名言警句来摆事实，讲道理，进行严密的论证，则与一般议论文的写法相同。这种文体和一般命题作文的其他文体的界限比较明显，因为它必须从“读”什么或“看”什么出发进行议论。它就是和一般的文章介绍、短评等也有所区别。它是根据原作的内容抒写感想，文章介绍则重在说明原文的内容和形式，它偏重于对主观感受的抒发，短评则偏重于对客观事物的评论。我们经常进行这种类型的写作训练，可以加深对原作的认识和理解，可以提高对原作的欣赏水平，更主要的是可以培养我们的分析能力、联想能力和表达能力。

那么，我们怎样来写读后感、观后感呢？

1. 精读原作或揣摩原作，认真分析原作的字面意义或画面意义，掌握原作的比喻意义、引伸意义，透过原作所表现的事件、人物或论点概括出原作的精神实质，从中引出自己感受最深的地方，确定自己的“感点”。这是写好读后感或观后感的关键一步。

大家知道，不管是写读后感，还是写观后感，首先得“读”或“观”，然后才有“感”。“读”或“观”是“感”的基础。只有在仔细的“读”或“观”的基础上，才可能领会原作的精神实质，并且才可能受到原作的感染和触动，以致选准“感点”，写出情真意切的感想。如对原作本身理解肤浅，或者舍本逐末，钻牛角尖，那么写出来的文章，不是空洞浮泛，缺乏新意，就是感受不深，议论不准，甚至会歪曲原意，而弄得南辕北辙。

例如，有这样一则短文：
\begin{center}
	改\hspace{1em}对\hspace{1em}联
\end{center}

明末东林党人顾宪成曾写过一副有名的对联：“风声、雨声、读书声、声声入耳；家事、国事、天下事，事事关心。”现在北京有位同志出于对“四人帮”造成的不良社会风气的忧愤，把顾宪成的对联改成“风声、雨声、不哎声，了此一生；家事、国事、不问事，平安无事。”一位老山战士看了以后，认为这太消极了，不符合我们这个时代的要求，又把这副对联改成“枪声、炮声、杀敌声，壮丽一生；苦事、难事、没有事，无畏战士。”

有位老师把这则短文作为材料发给大家，叫同学们以《读〈老山战士改对联〉有感》为题，写篇议论文。结果有的同学草草看了一下材料，就把“既要努力读书，又要关心政治”作为自己的“感点”，实际是把课文中马南邨（邓拓）同志所写的《事事关心》对于顾宪成这副对联的中心感受照搬了过来，既无新意，又不切合这则短文的精神实质和时代气息；还有的同学只取北京那位同志的改联作为触发“感点”的对象，把“歪风邪气泛滥成灾，不能不使人扼腕愤恨”作为中心感受，完全抛开了顾宪成的原联和老山战士的改联，议论不准也就是必然的了。更多的同学则融会贯通这则短文的含义，对比了三副对联的内容，确定“老山战士改对联”的“改”字是原作的“题眼”，因而把“感点”集中在老山战士针对什么“改”和怎样“改”上，从而选准“不能面对时代风雷枉自悲叹，必须挺身而出，知难而上，争当时代的勇士”作为自己的中心感受。这个感点选得准，选得好。准，是指准确地反映了这则短文的精神实质；好，是指正确地回答了人们普遍关心的问题，具有深刻而强烈的时代意义。

又如一九八三年全国高考语文试卷的作文题之一，是提供一幅《这下面没有水，再换个地方挖！》的漫画，叫考生根据这幅漫画的内容，自拟题目，写一篇议论文，实际上就是要求考生写一篇观后感。有些考生没有仔细观察画面上出现的人、物、景和它们之间的内在联系，因此就不能把握图画的寓意，作出正确的判断，于是说什么“越寇挖大坝没有好下场”呀，什么“必须警惕反革命分子的破坏和捣乱”呀，什么“挖不到石油的教训”呀，什么“单干没有出路”呀，什么“不改革不行”呀，五花八门的离奇感受都出来了。相反，很多考生首先仔细、全面地观察了画幅下部的水层，水层上面的土层，土层纵剖面的左侧五个深浅不一的干井（有两个已接近了水层），土层右上角一位口中含烟、右臂挎着衬衫、左手拿着铁锹、裤脚高高挽起、正跨步向前的小伙子，画幅之外的画题：《这下面没有水，再换个地方挖！》然后分析了这几者之间的联系，从而把握了画图的寓意，作出正确的判断，选准自己的感点，锲而不舍必能成功，万事都要深入开掘；为半途而废者戒，等等。这些感点也是选得准、选得好的。说它准，因为能准确地反映出这幅漫画的寓意；说它好，因为能给自已和读者以深刻的启示。

2. 开始下笔时，一般先把“读”或“观”的内容扼要写出，或直接引用原文中触及自己“感点”的某些语句，以突出生发感想的基础，从而引出自己所“感”的中心论点。这是一般读后感或观后感的开头法。

大家知道，“读”或“观”只是“感”的前提，而“感”才是文章的主体。因此，这个“感”的“前提”或“基础”，既要涉及原作的精华，又要表述简约、精当；决不可照抄原作，也不能详细复述原作的内容。

例如一九八〇年全国高考语文试卷中的作文题《读<画蛋>有感》，原作《画蛋》有三百多字。有的考生在文章开头几乎照抄原作，不会概述，也不会恰当引述，更没有突出触发感想的内容，这当然不符合要求。记得当年有位北京考生王延风，他写的开头真是别开生面。他这样写道：

人说“名师出高徒”，这话实是不假。依我看，这原因有二：其一，是师之教有方；其二，是徒之学不怠。离开了这两条，既称不上“名师”，也出不了“高徒”。名画家佛罗基奥手下出了一代宗师达·芬奇，便是很好一例。

你看，他一下笔就提出了问题，开门见山地突出中心论点：名师出高徒，在于“师之教有方”，“徒之学不怠”。
这是他在抓准原作的精神实质的基础上提炼出来的独特感受。而概述原作内容的部分，只有最后一句话，这句话恰恰又有力地支持了他的文章的中心论点。这样的开头，是何等简炼、何等鲜明、何等新颖！

再如一九八一年全国高考语文试卷的作文题《读〈毁树容易种树难〉》，原作本来只有百把字。有的考生开头一段复述原作的话就有二、三百字，还没触及“感点”；相反，有的考生只用两三句话就既概述了原作，又点明了中心感受。如“读罢《毁树容易种树难》，深深有感于毁树之易，种树之难。它教我悟出了这样一个道理：无论做什么事情，都是不容易的，但要破坏它却易如反掌。”它既突出了必须引述的内容，又为全文论说的中心感受打下了基础，这样的开头不是也很简洁、恰当吗？

又如写《看电影〈少年犯〉有感》，不少同学一下笔就沉溺在影片情节的复述中，写了三、四百字还意有未尽，似乎写得颇曲折、颇全面，但是引起你写这篇文章的“感点”是什么却十分含糊。有的同学就不然了，他只引述了几句主题歌词，很快就突出了文章的中心议题，鲜明地点出了“感点”，他是这样开头的：

“妈妈呀，儿已跌入急流，几番沉浮，不能自拔。又恰似狂风暴雨，摧折了未放的花。……妈妈呀，有妙手回春，残枝败叶又放新花，儿已被扶上骏马，去追回失去的年华。妈妈呀，妈妈呀，妈妈呀，妈妈呀······”我听着这撕裂人心的歌声，一边泪水禁不住哗哗直流，一边心头禁不住思潮翻滚：亲爱的老师、家长和社会上的人们啊，请伸出手，救救这些误入歧途的孩子吧！

你看，他多么恰当地引述了《少年犯》的主题歌词。既概述了《少年犯》的基本内容，又水到渠成地引发了自己的感受。这感受，既有形象的描写，又有抽象的概括，十分自然地突出了自己文章的中心论点。这种写法，不是很值得我们学习吗？

3. 就说明自己感受的中心论点，由小到大、由近到远、由此及彼、由人及已地进行类比，展开联想，举出典型事例或名言警句，进行由浅入深、由表及里的分析和论证。这是全文的主干部分。

大家知道，不论是写读后感还是写观后感，“读”或“观”只是“感”的基础，“读”或“观”在文章布局上应为辅，而“感”应为主。“感”得是否正确、深刻、新颖，是决定这类文章水平高低、议论深浅的标志。因此，在篇幅上，“感”要多于“读”或“观”，如果把原作内容从头到尾写出来，抓不住重点，在“感”字上写得过少，则将喧宾夺主，不合要求；在发挥感想上，一定要围绕“感点”，展开联想，联系实际，如果将“感”的内容只停留在“读”或“观”的内容上，不会抓住本质，展开联想，而只是就事论事，在原作的字面上兜圈子，那么这篇读后感或者观后感必定是肤浅的；在论据表述上，应该简明扼要，力求概述，如果将引用的事例写得太详，甚至成了有头有尾的记叙文，使议论部分显得又少又弱，也不符合作为议论文的读后感或观后感的要求；在抒发感情上，必须是真情实感，是自己感受最深的一点，如果纠缠于旁枝末节，甚至无病呻吟，而不是抒写自己的真情实感，没有突出自己感受最深的一点，当然也是不合要求的；在联系实际上，应该恰当、准确，有针对性，有教育意义，选择有代表性的思想或事例进行实事求是的分析，如果只是简单地不分轻重、不分主次地罗列事实，或者把联系实际写成检讨书或决心书，就都与有感而发、切中要点背道而驰了。

一般说来，紧扣“感点”，进行恰当的类比联想，是发挥感想、论证中心的好办法。这种联想，可以是“感点”的生发，也可以是事实的类比或扩大。例如写《读<老山战士改对联>有感》，就老山战士的改联生发出“不能面对时代风雷枉自悲叹，而要争当无畏的战士”，然后以一些敢于与不正之风斗争者，敢于冲破阻力进行改革者或敢于为实现四化英勇献身者的典型事例进行论证，就是对“感点”的生发；又如写《读<毁树容易种树难>》，想到“毁林容易育林难”，则是简单类比；想到“毁才容易育才难”，则是进层类比；甚至想到“败坏党风容易，恢复党风艰难”，“毁事容易成事难”等等，则把类比的范围更加扩大了。当然，就事论事，简单类比，感想还比较肤浅，能够进层类比，扩大类比的范围，只要紧扣“感点”，类比恰当，就能使感想发挥得淋漓尽致，中心感受也就显得越加深刻了。

我们可以用北京那位考生王延凤写《读<画蛋>有感》为例，看看怎样才能使自己的感受发挥得淋漓尽致，从而对中心论点进行逻辑严密的论证。他先紧扣《画蛋》原文中心即要苦练基本功，提出自己的中心论点：“师之教有方，徒之学不怠。”是“名师出高徒”的根本原因。然后围绕这个中心“感点”，展开联想，联系实际，举出了两个与原文相类似的典型事例：一是自己向一位女老师学拉小提琴，她要自己先要苦练拉空弦这个基本功，而自己最后一烦之下，竟然不学了，这就与原文中所写的“教有方”“学不怠”进行了对比，以自身经历写出只是“师之教有方”而“徒之学怠”，仍然不能成功；二是自己向一位老师学写作，老师要自己天天记札记，自己坚持下来了，终于在写作上有了进步，这就与原文进行映照，同样以自身的经历作为例证，进一步论述了“师之教有方”加上“徒之学不怠”而取得成功的道理。这两个类比联想一反一正，层层深入，对中心感受作了深刻有力的论证，能使读者从中得到启示和教益。

写读后感如此，写观后感也应如此。我们只要善于紧扣“感点”，展开联想，进行有力的分析和论证，就一定能够“感”得深，想得透，写得好。

4. 结尾时，一般要照应原作或照应引述的原作中的要点或警句，使“读”或“观”能与“感”有机结合，从而概括中心，总结全文，或者提出问题，深化主题，发人深思。

这个部分尽管篇幅不长，但它有关全篇结构的完整、布局的谨严、层次的分明、条理的清楚，特别是有关中心的深化或强化，切切不可用几句空话套语敷衍了事。我们要使文章有“凤头豹尾”之姿，而不可有“虎头蛇尾”之态。人们常有这样的感觉：吃最后一颗花生米喷香满口，能让人长久回味；如果最后一颗花生米是焦糊恶苦的，那么前面的香味将被败坏无余了。

只要我们认真地对待全文的收尾，就会精心构思尾段的语意了。这里至少应该注意如下几点：

(1)巧妙呼应原作或所引用的原作的要点、警句，使人感到“感”从“读”或“观”而来，“读”或“观”又为“感”所用。

(2)灵活地照应开头或前文，使文章脉络贯通，前呼后应。

(3)画龙点睛，强化中心感受，用语简洁精警。

(4)突发奇想，深化中心感受，使人产生言有尽而意无穷之感。当然，如果能把这几点综合运用起来，就更好了。

例如北京考生王延风的《读〈画蛋〉有感》是这样结束的：

芬奇画蛋直至学画成功和我的失败与进步，都说明不仅教者要有方，学者更应不怠，这两方面的因素结合起来，能有什么知识学不会，有什么事情做不好呢？他不仅呼应了原文，照应了上文，更用反问的句式有力地强化了中心感受，成了全文的画龙点睛之笔。

再如一位南京考生在写《读〈毁树容易种树难〉》时，以“毁才容易育才难”为中心感受，然后他以“纵观古今中外，多少非凡的人才，逆境中成长，逆境中成才，显示出非凡的生命力，然而他们能幸免于‘天灾’，却往往夭折于‘人祸’！”和“毁树之易与毁人之易，都是史有前例的；然而，在十年浩劫时期，毁树、毁人却容易到史无前例的地步！”两大层意思，围绕“感点”，展开联想，列举古今典型事例，深刻有力地论证了“毁才容易育才难”这个中心感受。最后，他这样结尾：

值得欣慰的是，这噩梦般的年代已经过去了。我望着窗外的法国梧桐，想：种树诚然难，在今天毁树又谈何容易呢？

然而，我们竟可以放松警惕吗？

这真是奇峰突起，出人意料，但是细想之下，说得又何等在理，何等贴切。它不仅紧扣原文，呼应前文，强化中心，而且思路一波三折，引人入胜，联系现实具体、恰当，立意新颖精警，有力地深化了主题。这样认真深入地思考结尾的写法，不是很值得我们学习吗？

总起来说，读后感或观后感没有固定格式，写法比较灵活，我们所说的只是一般的写法。它往往采用夹叙夹议的方式进行表达，以“议”为主，以“叙”为副。在“读”、“观”同“感”的关系上，“读”、“观”是基础，“感”是主体，“感”从“读”“观”而来，“读”“观”为“感”所用。文章的开头一般需要点破“读后”或“观后”，即你读什么文章，中心思想是什么，它给人什么启示，这类文字要少而精，不要拖沓；有时可以转述或摘录原作中触发“感点”的句子，适当阐发，以确定中心感受；文章的主体部分，应该扣住原作，展开类比联想，紧密联系实际，进行恰当发挥，做到有破有立，说理透辟，结束语要总结全文，照应开头，做到简短有力，首尾圆合。其中最重要之点，是吃透原文，选准感点，务求立意正确、鲜明、深刻和新颖。我们只要掌握这些基本的常识和习用的方法，再揣摩范文，认真练习，完全可以写好读后感或者观后感。
\newpage